% !TeX root = ../main.tex
\section{Indeks Bukti dan Jejak Eksekusi}
\label{app:evidence}
Seluruh berkas E1 hingga E8 berada di \nolinkurl{docs/evidence/integration/} pada revisi acuan. Artefak disimpan sebagai teks dan JSON agar dapat ditinjau tanpa bergantung pada screenshot. Potongan yang dicetak merupakan petikan atau rekap berlabel; output lengkap tetap menjadi acuan.

\begin{tabularx}{\linewidth}{lP{0.38\linewidth}Y}
\toprule ID & Berkas & Isi\\\midrule
E1 & \evidence{modules.txt} & Uji unit dan go vet sembilan module dan bootstrap.\\
E2 & \evidence{regression.txt} & Regresi lengkap termasuk auth, ingest, query, drift, dan rebuild.\\
E3 & \evidence{events-fresh-group.txt} & Replay subscriber baru, DLQ, dedup, dan outage broker.\\
E4 & \evidence{load/seismic-only.json}, \evidence{load/sustained.json}, \evidence{load/outage.json} & Summary k6 awal. Angka akhir laporan memakai E9.\\
E5 & \evidence{load/connections.json} & Sampel waktu host dan koneksi TCP teramati.\\
E6 & \evidence{event-trace.jsonl} & Trace nyata dari mock sampai consumer.\\
E7 & \evidence{postgres.txt} & Integrasi transaksi dan query PostgreSQL.\\
E8 & \evidence{services.txt} & Snapshot status container akhir sesi pengujian.\\
\bottomrule
\end{tabularx}

\subsection{Bukti pengujian akhir}
E9 merujuk \finalevidence{README.md} beserta JSON k6, sampel koneksi, dan konfigurasi pada subfolder load. Angka P2 memakai pengujian final ini.

E10 merujuk \finalevidence{regression.txt} dan \finalevidence{postgres.txt}. Regresi stack lulus dalam 306,263 s. Bukti unit seluruh module dan scan histori yang cakupannya lebih awal tetap tersedia pada \reliabilityevidence{modules.txt} serta \reliabilityevidence{secret-scan.txt}.

E11 merujuk \codepath{docs/evidence/http-attempts-2026-10-08/README.md}, yang memuat hasil tes module Client API dan go vet untuk verifikasi latency per percobaan HTTP. Bukti ini tidak menggantikan pengukuran k6 pada E9.

E12 merujuk \finalevidence{long-outage/result.json}, \finalevidence{long-outage/summary.json}, dan \finalevidence{bootstrap/result.json}. Artefak ini memeriksa outage 600 detik, pemulihan tanpa restart, serta bootstrap dengan kredensial dan volume baru. Manifest pada \finalevidence{manifest.json} mencatat hash artefak serta runner.

\subsection{Contoh trace event nyata}
Correlation ID: \nolinkurl{2a877ebd-3298-4ba1-8f24-8bed0d80d6e3}.\par\smallskip
\begin{tabularx}{\linewidth}{lYr}\toprule
Service & Pesan log & Latency (ms)\\\midrule
bmkg-mock & http\_request & 100.426\\
aggregator & poll\_complete & 118\\
aggregator & outbox\_published & 1\\
pemda-portal & record\_completed & 6\\
notifier & notification\_simulated & --\\
dashboard-updater & record\_completed & 5\\
notifier & record\_completed & 4\\
\bottomrule\end{tabularx}

Correlation ID dibentuk pada polling lalu diteruskan ke mock, transaksi dan envelope, publish, dan consumer. Query client memiliki correlation ID tersendiri; query tidak dianggap sebagai penyebab ingest. Pada log consumer, latency menunjukkan waktu pemrosesan record sampai commit, bukan latency ujung-ke-ujung sejak kejadian bencana.

\subsection{Asal potongan kode dan gambar}
Potongan kode di \nolinkurl{assets/code/} diekstrak dari revisi acuan, dengan lokasi dan hash pada \nolinkurl{assets/provenance.json}. Sumber enam diagram berada di \nolinkurl{diagrams/}; grafik performa dibangkitkan dari JSON bukti final. Semua dapat dibuat ulang dengan \nolinkurl{scripts/refresh-assets.py}. Gambar bukan hasil tangkapan layar implementasi atau bukti tambahan yang dibuat-buat.
